\section{Тема 14. Плоскость в трёхмерном пространстве}

\subsection{Постановка задачи (вариант~4)}

Плоскость в пространстве задана общим уравнением
$Ax+By+Cz+D=0$. Найти расстояние от заданной точки $(x_0,y_0,z_0)$
до этой плоскости и координаты проекции точки на плоскость.

\subsection{Математическое обоснование}

Расстояние от точки до плоскости вычисляется по формуле:
\[
  \rho = \frac{|Ax_0+By_0+Cz_0+D|}{\sqrt{A^2+B^2+C^2}}.
\]
Проекция точки на плоскость находится движением вдоль вектора
нормали $\vec n=(A,B,C)$ на параметр
\[
  t = -\frac{Ax_0+By_0+Cz_0+D}{A^2+B^2+C^2},
\]
так что координаты проекции:
\[
  (x_0+At,\;y_0+Bt,\;z_0+Ct).
\]

\subsection{Блок-схема алгоритма}

\begin{center}
\begin{tikzpicture}[
  node distance=10mm, start chain=going below,
  nodes={draw=black, top color=white, bottom color=lightgray!50,
         minimum width=5.6cm, align=center, font=\small},
  arrow/.style={->, >=Stealth, thick},
]
  \node[draw, rounded corners, on chain] {Start};
  \node[shape=trapezium, trapezium left angle=70,
        trapezium right angle=110, on chain, draw, join=by arrow]
       {$A,B,C,D$;\ $x_0,y_0,z_0$};
  \node[shape=rectangle, on chain, draw, join=by arrow]
       {$\|n\| := \sqrt{A^2{+}B^2{+}C^2}$};
  \node[shape=rectangle, on chain, draw, join=by arrow]
       {$t := -(Ax_0{+}By_0{+}Cz_0{+}D)/\|n\|^2$};
  \node[shape=rectangle, on chain, draw, join=by arrow]
       {$\rho := |Ax_0{+}By_0{+}Cz_0{+}D|/\|n\|$};
  \node[shape=rectangle, on chain, draw, join=by arrow]
       {$Q := (x_0{+}At,\ y_0{+}Bt,\ z_0{+}Ct)$};
  \node[shape=trapezium, trapezium left angle=70,
        trapezium right angle=110, on chain, draw, join=by arrow]
       {$\rho$,\ $Q$};
  \node[draw, rounded corners, on chain, join=by arrow] {End};
\end{tikzpicture}
\end{center}

\subsection{Текст программы}

\lstinputlisting[style=Cstyle]{../src/t14.c}

\subsection{Тестовые данные}

\begin{center}
\begin{tabular}{l|l}
\toprule
Плоскость, точка & Расстояние, проекция \\
\midrule
$z=0$, $(1,2,5)$ & $\rho=5$, проекция $(1,2,0)$ \\
\bottomrule
\end{tabular}
\end{center}

\textbf{Результат выполнения программы:}
\begin{lstlisting}[language={},basicstyle=\ttfamily\small,frame=single]
Коэффициенты плоскости (A B C D): 0 0 1 0
Точка (x0 y0 z0): 1 2 5
Расстояние от точки до плоскости: 5.0000
Проекция точки на плоскость: (1.0000, 2.0000, 0.0000)
\end{lstlisting}
